\section{8.2 Peer-Review}
\begin{frame}{왜 Peer-Review인가: 이번 주에는 네가 심사위원이다}
  \begin{itemize}
    \item 4주 파이프라인에서 4주째는 \textbf{역전} --- 발표팀이었던 각
      팀이 \textbf{다른 두 팀의 발표를 심사}하고 리뷰 \textbf{2통}을 작성.
    \item 채점 대상은 리뷰의 \emph{양}이 아니라 \emph{질}: ``4항목을
      다 썼는가''가 아니라 \textbf{``그 질문이 반박 가능한 질문인가''}.
  \end{itemize}
  \vskip0.6em
  \begin{block}{근거 (ML1 8.2에서 그대로 이어받음)}
    \begin{itemize}
      \item Popper의 \textbf{반증주의}: 한 주장은 ``그것이 틀렸음을
        보여주는 관측이 상상이 가능한가''로 판별.
      \item Amgen 재현 실험(2011): 53편의 'landmark' 논문 중 \textbf{6편
        만} 재현. \kq{좋은 결과는 스스로를 증명하지 못하고, 결과가
        남이 재현·반박할 수 있도록 만들어졌는지가 기준.}
      \item 강화학습에서는 이 위험이 \textbf{더 높음} --- 결과를 내는
        \textbf{과정 자체}가 무작위(시드, 탐험)라 ``우연히 좋은''
        확률이 학습모델보다 훨씬 큼.
    \end{itemize}
  \end{block}
\end{frame}
